\documentclass[activite]{thedocument}
\startdatakeys
    \source{originale}
    \BO{2026}
    \cycle{4}
    \competencesDuSocle{CA, RE}
    \objectifsApprentissage{C4A1cb, C4A1ca}
\enddatakeys
\begin{document}
    \begin{enumerate}
        \item Pour son anniversaire, les amis de Marie
        lui ont préparé un gateau. Enzo a mangé un cinquième
        du gateau et Marie en a mangé deux cinquièmes.\par
        \vskip5pt\noindent
        Quelle fraction du gateau Marie et Enzo ont-ils mangé
        a eux deux~?
        \answerspace{30pt}
        \item En déduire une règle permettant d'additionner deux
        fractions ayant le même dénominateur.
        \answerspace{50pt}
        \item Un autre gâteau est préparé pendant la fête. Estelle
        en mange deux neuvièmes et Nimra trois fois plus.\par
        \vskip5pt\noindent
        Quelle fraction du gâteau Nimra a-t-elle mangé~?
        \answerspace{30pt}
        \item En déduire une règle permettant de multiplier une
        fraction par un nombre entier.
        \answerspace{50pt}
    \end{enumerate}
\end{document}